% ============================================================================
\chapter{Cuando algo cambia}\label{cap:cambios}
% ============================================================================

La regla general: \textbf{lo que cambia desde una fecha no debe alterar el pasado}. Cada caso tiene su forma correcta.

\section{Cambia un precio o una tarifa}
Diésel, tarifa por kg, cita, dieta, pago por jornal, servicios de transitaria, comisión\dots
\begin{pasos}
\paso En \hoja{PRECIOS}, copie la última fila con precios en la fila siguiente vacía.
\paso Cambie la fecha de la columna A por la fecha desde la que rige el cambio.
\paso Cambie solo el precio que varió y anote el documento en la columna O.
\end{pasos}
Es el ejercicio~4 (sección~\ref{sec:ej4}). Mantenga las fechas \textbf{en orden}, de la más antigua a la más reciente.

\section{Cambia un gasto fijo de forma permanente}
Por ejemplo, sube el salario de la oficina desde marzo.
\begin{pasos}
\paso En \hoja{COSTOS\_FIJOS}, vaya a la fila del concepto y a la columna de \textbf{marzo}.
\paso Escriba el nuevo importe. Cópielo con \tecla{Ctrl}+\tecla{C}, seleccione las celdas de los meses siguientes y péguelo con \tecla{Ctrl}+\tecla{V}.
\paso \textbf{No toque el valor base} (columna F): así los meses anteriores conservan el importe antiguo.
\end{pasos}

\section{Cambia una norma técnica}
Si después de varios contenedores comprueba que una norma ya no es real (por ejemplo, el camión consume 70 L y no 60), cámbiela en \hoja{PARAMETROS}.
La norma solo afecta a las celdas opcionales que estén \textbf{vacías}. Por eso conviene escribir siempre los litros reales: así un cambio de norma
nunca altera el pasado.

\section{Se compra, se vende o se retira un vehículo o equipo}
\begin{itemize}
  \item \textbf{Compra}: agréguelo en una fila libre de \hoja{ACTIVOS} con su fecha de alta. La depreciación empieza ese mes.
  \item \textbf{Venta o retiro}: escriba la \textbf{fecha de baja} en la columna K. Desde el mes siguiente deja de depreciarse y de cargar mantenimiento.
        \textbf{No borre la fila}: los meses anteriores la necesitan.
\end{itemize}

\section{Cambia un gestor}
\begin{itemize}
  \item \textbf{Sustitución} en el mismo territorio: cambie el nombre en la columna B de \hoja{COMERCIAL}.
  \item \textbf{Cambio de pago fijo para todos}: para que no cambien los meses anteriores, siga la sección «Cambia un gasto fijo de forma
        permanente»: escriba el nuevo total de la fila 8 de \hoja{COSTOS\_FIJOS} en los meses desde el cambio. Luego actualice \celda{C8} de \hoja{COMERCIAL}.
\end{itemize}

\section{Cambian los impuestos}
Actualice las tasas en \hoja{PARAMETROS} (sección D) \textbf{después de confirmarlo con el contador}.

\begin{atencion}
Las tasas de \hoja{PARAMETROS} se aplican a \textbf{todos} los meses del libro. Si cambian a mitad de año, guarde antes una copia del libro con
las tasas antiguas como histórico, o anote el cambio para tenerlo en cuenta en la liquidación anual.
\end{atencion}

\section{Se llena el libro: nuevo período}
El libro admite \textbf{24 meses} y \textbf{300 contenedores}. Al acercarse a cualquiera de los dos límites:
\begin{pasos}
\paso Haga un respaldo y guarde el archivo como histórico, por ejemplo \escribal{Rentabilidad\_\allowbreak HISTORICO\_\allowbreak 2026-2027.xlsx}.
\paso Haga una copia como nuevo maestro.
\paso En el nuevo maestro, cambie \texttt{Fecha\_Inicio} en \hoja{PARAMETROS} al primer mes del nuevo período.
\paso Borre los contenedores de \hoja{ENTRADA} (con \tecla{Supr}) y revise que \hoja{PRECIOS}, \hoja{ACTIVOS} y \hoja{COSTOS\_FIJOS}
      tengan los valores vigentes.
\end{pasos}

\section{Ver o cambiar algo protegido}
Las hojas están protegidas \textbf{sin contraseña}, solo para evitar errores. Si necesita cambiar algo bloqueado:
\menu{Revisar $\rightarrow$ Desproteger hoja}. Haga el cambio y vuelva a protegerla con \menu{Revisar $\rightarrow$ Proteger hoja}
$\rightarrow$ \boton{Aceptar}, sin escribir contraseña.

\begin{atencion}
Con la hoja desprotegida, las fórmulas se pueden borrar sin querer. Si algo falla después de un cambio así, recupere la última copia de respaldo.
\end{atencion}
